% Software Impacts usa la clase cas-dc (doble columna) del paquete
% oficial de Elsevier "els-cas-templates" -- NO elsarticle. La
% estructura de contenido (Code metadata con exactamente los campos
% C1-C9, cuerpo libre de hasta 3 paginas, Acknowledgements opcional,
% References) sigue al pie de la letra la plantilla oficial:
% https://legacyfileshare.elsevier.com/promis_misc/SIMPAC_Article_Template.pdf

\documentclass[a4paper,fleqn]{cas-dc}

\usepackage[numbers]{natbib}
\usepackage{booktabs}
\usepackage{enumitem}

\begin{document}

\let\WriteBookmarks\relax
\def\floatpagepagefraction{1}
\def\textpagefraction{.001}

\shorttitle{Campus-DDoS-SDN}
\shortauthors{Calder\'on Leiva et al.}

% Titulo propuesto, siguiendo la misma convencion de nombres que usan
% Kumar y Gupta en sus dos datasets de referencia ("SDN TCP-SYN
% Dataset", "ICMP-Flood-SDN"): <nombre-corto>: descripcion. Cambiar el
% nombre corto si el equipo prefiere otro.
\title[mode=title]{Campus-DDoS-SDN: A Python-based testbed for
multi-vector DDoS attack detection in software-defined campus networks}

\author[1]{Ana Olenka Calder\'on Leiva}

\author[1]{Ricardo Llanos Lozano}

\author[1]{Cristhian Emilio Mart\'inez Casas}
\cormark[1]
\ead{cristhianmartinezcasas@gmail.com}

\author[1]{Alessandro Piero Poves Mart\'inez}

\cortext[1]{Corresponding author}

\affiliation[1]{organization={Universidad Nacional del Centro del Per\'u},
            addressline={},
            city={Huancayo},
            postcode={},
            state={},
            country={Peru}}

\begin{abstract}
% Software Impacts pide ~100 palabras. TODO: completar una vez esten
% los resultados de la Fase 5.
TODO: resumen de ~100 palabras. Debe cubrir: que hace el software
(entorno reproducible Mininet+Ryu+ML para generar y clasificar
trafico DDoS multi-clase), como se diferencia del trabajo de
referencia (deteccion de 4 tipos de ataque en vez de solo ICMP flood,
entorno reproducible via Docker), y el resultado principal obtenido.
\end{abstract}

\begin{keywords}
Software defined networking \sep
Machine learning \sep
DDoS attack detection \sep
Mininet \sep
Network traffic classification
\end{keywords}

\maketitle

%% ------------------------------------------------------------------
%% Tabla de metadata del codigo: EXACTAMENTE los campos C1-C9 que pide
%% la plantilla oficial de Software Impacts (no inventar columnas
%% propias).
%% ------------------------------------------------------------------
\begin{table*}[t]
\small
\caption{Code metadata.}\label{tbl:metadata}
\begin{tabular*}{\textwidth}{@{}p{1cm}p{6.5cm}L@{}}
\toprule
Nr. & Code metadata description & Please fill in this column \\
\midrule
C1 & Current code version & v1.0 \\
C2 & Permanent link to code/repository used for this code version & TODO: URL del repo \\
C3 & Permanent link to Reproducible Capsule & -- \\
C4 & Legal Code License & TODO: elegir licencia (p.ej. MIT) \\
C5 & Code versioning system used & git \\
C6 & Software code languages, tools, and services used & Python, Bash, Mininet, Open vSwitch, Ryu, CICFlowMeter, Docker \\
C7 & Compilation requirements, operating environments \& dependencies & WSL2/Ubuntu 24.04 o Docker; Python 3.9 (Ryu) y Python 3.12 (CICFlowMeter); ver \texttt{README.md} \\
C8 & Link to developer documentation/manual & Ver \texttt{README.md} del repositorio \\
C9 & Support email for questions & TODO: email de contacto \\
\bottomrule
\end{tabular*}
\end{table*}

%% ------------------------------------------------------------------
%% Cuerpo del articulo (hasta 3 paginas, excluyendo metadata, tablas,
%% figuras y referencias). La plantilla no exige subtitulos fijos,
%% pero pide cubrir estos puntos: funcionalidad de alto nivel para
%% audiencia no especialista, "impact overview", proyectos en curso,
%% limitaciones/trabajo futuro, y publicaciones habilitadas por el
%% software.
%% ------------------------------------------------------------------

\section{Introduction}\label{sec:intro}

Software-defined networking (SDN) decouples the control plane from
the data plane, concentrating network intelligence in a logically
centralized controller \citep{kreutz2015sdn}. This centralization
simplifies network management, but it also turns the controller into
a single point of failure: a distributed denial-of-service (DDoS)
attack that saturates the controller, the switches, or a hosted
service can degrade or take down the entire network segment it
governs. Detecting such attacks early, and separating them reliably
from legitimate traffic spikes, requires datasets that are both
representative of real SDN deployments and precisely labelled.

Most publicly available intrusion-detection datasets, such as
NSL-KDD \citep{tavallaee2009nslkdd}, CICIDS2017
\citep{sharafaldin2018cicflowmeter}, and CICDDoS2019
\citep{sharafaldin2019cicddos2019}, were not captured inside an SDN
environment: they lack OpenFlow-specific signals, such as flow-table
occupancy or controller load, and are nonetheless often reused out of
context to train SDN-targeted classifiers. Recent work has started to
close this gap by generating DDoS datasets directly inside Mininet/Ryu
testbeds. Kumar and Gupta released two single-vector datasets built
this way: a TCP-SYN flood dataset \citep{kumar2025synflood} and, more
recently, an ICMP-flood dataset \citep{kumar2026icmpfloodsdn} in which
normal traffic is captured by connecting a Mininet host to a live
internet connection and browsing real websites (YouTube, Facebook,
news, etc.), instead of relying on synthetic traffic generators. This
real-traffic approach increases realism, but it introduces two
practical limitations for a reusable benchmark: the resulting dataset
is tied to a single attack vector, and it is not reproducible across
runs, since the captured content and timing depend on whatever the
live internet happens to serve at capture time.

This work, Campus-DDoS-SDN, addresses both limitations within a
single, fully reproducible testbed. Four attack vectors (TCP-SYN
flood, HTTP flood, Slowloris, and ICMP flood) are launched from
already-compromised hosts against one shared victim, a web portal,
inside a Mininet/Ryu topology that mimics a university campus: a
core/DMZ switch layout plus one emulated laboratory per academic
area. Unlike prior single-vector or fully-separated-capture datasets,
legitimate users keep browsing the portal throughout every attack, so
each capture mixes benign and malicious flows against the same
target, as a real DDoS incident would. Ground-truth labels are
assigned per source IP rather than per capture file, and every run
additionally logs the attack's effect on the victim as perceived by a
real client (HTTP response code, latency, active connections),
alongside network-level indicators (packets and bytes per second,
OpenFlow flow-table size). The complete environment, including the
traffic generators, the custom web portal, and the labelling
pipeline, is packaged as a Docker image to remove any dependency on a
live internet connection or on a specific host operating system.

The main contributions of this work are:
\begin{enumerate}[label=\roman*.]
  \item A multi-class DDoS dataset covering four attack vectors
        against a single victim inside a realistic, multi-segment SDN
        topology, instead of a single attack type against a flat
        network.
  \item Per-flow, per-source-IP labelling that keeps legitimate and
        malicious traffic mixed within the same capture, closer to a
        real attack scenario than the fully-separated captures used
        in prior work.
  \item Combined network-level (flow and packet statistics via
        CICFlowMeter) and application-level (HTTP status, latency,
        connection saturation) impact metrics, captured automatically
        on every run.
  \item A one-command, fully reproducible Docker environment that
        does not depend on a live internet connection, addressing the
        reproducibility limitation of real-traffic-based SDN
        datasets.
\end{enumerate}

\section{Software description}\label{sec:description}

\subsection{Architecture overview}
Topologia Mininet de 5 hosts (servidor propio, 2 clientes, 2
atacantes) sobre un switch Open vSwitch (datapath userspace)
controlado por Ryu; generacion de 4 tipos de ataque (SYN flood, HTTP
flood, Slowloris, ICMP flood) capturados por separado; conversion a
CSV con CICFlowMeter (y un agrupador de flujos propio para ICMP, que
CICFlowMeter no soporta); todo empaquetado en una imagen Docker para
reproducibilidad completa. TODO: ampliar.

\subsection{Software functionalities}
TODO: enumerar los scripts principales (\texttt{scripts/mininet},
\texttt{scripts/wsl}, \texttt{scripts/docker}) y que hace cada uno.

\section{Impact overview}\label{sec:impact}
% Seccion explicitamente pedida por la plantilla de Software Impacts.
% Cubrir: que nuevas preguntas de investigacion habilita, en que
% forma/medida mejora la investigacion existente, si cambio la
% practica diaria de sus usuarios, que tan difundido esta su uso, y si
% se usa en entornos comerciales.
TODO: completar. Punto de partida: a diferencia de
\cite{kumar2026icmpfloodsdn}, que depende de una conexion a internet
real y por lo tanto no es reproducible entre corridas, este trabajo
habilita comparaciones reproducibles de metodos de deteccion
multi-clase de DDoS en SDN.

\subsection{Ongoing research projects using the software}
TODO (si aplica).

\subsection{Limitations and future improvements}
El trafico normal se genera de forma sintetica y controlada (en vez
de navegacion real a internet, como en \cite{kumar2026icmpfloodsdn}),
lo que reduce la diversidad/realismo del trafico a cambio de
reproducibilidad total. Trabajo futuro: deteccion en tiempo real
mediante una aplicacion Ryu propia que clasifique estadisticas
OpenFlow en vivo.

\subsection{List of scholarly publications enabled by this software}
TODO (si aplica, una vez publicado).

\section{Software and hardware requirements}\label{sec:requirements}

\begin{table*}[t]
\small
\caption{Testbed setup requirements.}\label{tbl:testbed}
\begin{tabular*}{\textwidth}{@{}p{1cm}LLp{2cm}@{}}
\toprule
N.\textsuperscript{o} & Descripcion & Herramienta & Version \\
\midrule
1 & Emulador de red & Mininet & TODO \\
2 & Controlador SDN & Ryu & 4.34 \\
3 & SYN / ICMP flood & hping3 & TODO \\
4 & Ataque de baja tasa & slowhttptest & TODO \\
5 & pcap a CSV & CICFlowMeter & 0.5.0 \\
\bottomrule
\end{tabular*}
\end{table*}

\section*{Acknowledgements}
TODO (opcional).

\bibliographystyle{cas-model2-names}
\bibliography{references}

\end{document}
